\documentclass[12pt]{article}

\usepackage[margin=1in]{geometry}
\usepackage{amsmath, amssymb}
\usepackage{booktabs}
\usepackage{array}
\usepackage{hyperref}
\usepackage[numbers,sort&compress]{natbib}
\usepackage{graphicx}
\usepackage{setspace}
\usepackage{xcolor}
\usepackage{caption}
\usepackage{lscape}
\usepackage{longtable}
\usepackage{subcaption}
\usepackage{lineno}
\usepackage{authblk}

\onehalfspacing

\title{\textbf{Estimating Infectious Disease Importation Risk
during the 2026 FIFA World Cup}}

\author[1]{Jose L. Herrera-Diestra}
\author[2]{Kaiming Bi}
\author[4]{Zeynep Ertem}
\author[5]{Amera Al-Amery}
\author[3]{Mallory Harris}

\affil[1]{Tarleton State University, Stephenville, TX, USA}
\affil[2]{UTHealth Houston, Houston, TX, USA}
\affil[3]{University of Texas at Austin, Austin, TX, USA}
\affil[4]{Binghamton University, Binghamton, NY, USA}
\affil[5]{Jordan University of Science and Technology, Irbid, Jordan}

\date{\today}

\begin{document}

\maketitle
\linenumbers

% ============================================================
\begin{abstract}
\noindent\textbf{Background.}
The 2026 FIFA World Cup will bring an estimated 1--5~million
international visitors to 11~US host cities between June~11 and
July~19, 2026---the largest tournament in history. Large-scale
international gatherings accelerate importation of infectious
diseases from diverse source populations. Advance estimation of
importation risk is essential for public health preparedness and
surveillance prioritization.

\medskip
\noindent\textbf{Methods.}
We developed a Poisson importation framework applied to five
diseases (dengue fever, influenza, malaria, measles, and pertussis)
across the 11~US venue cities. Three nested travel models of
increasing resolution were constructed: a baseline model using
routine June~2024 arrival data; a World Cup--adjusted model
incorporating projected visitor growth factors; and a
schedule-driven model routing WC fans to specific cities based on
match assignments. WHO incidence and BTS T-100 routing fractions were
combined with Monte Carlo uncertainty propagation (5,000 Uniform
draws on under-reporting and travel-while-infectious parameters)
to yield median importation estimates with 95\% uncertainty
intervals.

\medskip
\noindent\textbf{Results.}
Dengue posed the highest importation risk at most venue cities under
the schedule-driven model (median $\Lambda > 10$ expected importations
from Brazil alone; 95\% uncertainty interval 5.9--33.1), robust across the full literature-supported parameter range;
Atlanta was the exception, where malaria probability exceeded dengue,
driven by direct travel from West and Central African nations.
Influenza ranked second at most cities, coinciding with the Southern
Hemisphere winter peak. Pertussis showed broad geographic spread but
carries the widest relative uncertainty, as the assumed detection rate
sits at the upper bound of the literature range. Background tourism
accounted for the dominant share of total importation risk; the World
Cup fan increment contributed approximately 8.3\% of projected arrivals
for WC-qualified nations.

\medskip
\noindent\textbf{Conclusions.}
This Poisson importation framework, built entirely from publicly
available data, provides reproducible importation risk estimates for
mass gathering events. The framework extends to additional
diseases, cities, and gatherings, offering
a baseline complementary to proprietary modeling systems.
\end{abstract}

% ============================================================
\section*{Introduction}
% ============================================================

Mass gathering events concentrate large numbers of international
visitors in a short period, creating conditions for accelerated
importation of infectious diseases from diverse source populations.
The 2026 FIFA World Cup---the largest in the tournament's history,
with 48 qualified nations---will draw visitors from every world
region to 11 US host cities during June and July 2026, a period
that coincides with peak dengue transmission in the Southern
Hemisphere and with active global circulation of several
vaccine-preventable diseases.

Estimating importation risk in advance is essential for surveillance
prioritization and public health preparedness. Existing frameworks
often rely on aggregate passenger volumes that obscure spatial and
source-country heterogeneity, focus on a single disease, or do not
account for the event-specific travel surge. Here we present a
general, data-driven Poisson importation framework that addresses
these gaps, applied to the 2026 FIFA World Cup as a case study.

The framework integrates multiple publicly available datasets---US
Customs and Border Protection monthly arrivals (COR~I-94), BTS T-100
airport routing data, NTTO travel projections, and WHO disease
surveillance---into three nested models of increasing resolution. The
three-model structure allows clean attribution: the step from Model~1
to Model~2 isolates the effect of the World Cup travel surge; the
step from Model~2 to Model~3 reveals the added value of
schedule-based fan routing. The approach is reproducible using only
public data, and designed to be extensible to additional diseases and
geographies.

% ============================================================
\section*{Methods}
% ============================================================

\subsection*{Importation framework}

We model importation risk as a Poisson process. The expected number
of infectious travelers arriving at US venue city $v$ from country
$c$ for disease $d$ is:
\begin{equation}
  \lambda_{c,v,d} = N_{c,v} \cdot I_{c,d} \cdot p_d
  \label{eq:lambda}
\end{equation}
where $N_{c,v}$ is the projected number of travelers from country $c$
to city $v$, $I_{c,d} = \rho_d \cdot C_{c,d}/P_c$ is the
country-level disease incidence adjusted for under-reporting (with
$C_{c,d}$ reported cases, $P_c$ population, and $\rho_d$ the
under-reporting factor), and $p_d$ is the probability of traveling
while infectious. Total importation intensity at each venue city is
$\Lambda_{v,d} = \sum_c \lambda_{c,v,d}$, and the probability of at
least one importation is $P(\geq 1) = 1 - e^{-\Lambda_{v,d}}$.

\subsection*{Three-model hierarchy}

We applied the framework under three travel models of increasing
resolution (Table~\ref{tab:models}). All three models use the same
BTS T-100 routing fractions $f_{c,v}$ (mean June 2023--2025) to
distribute country-level arrivals across the 11 US venue cities, and
the same country-level disease incidence estimates.

\begin{table}[h!]
\centering
\small
\begin{tabular}{p{2.8cm} p{4.2cm} p{6cm}}
\toprule
\textbf{Model} & \textbf{Travel volume} & \textbf{Key feature} \\
\midrule
Model 1: Baseline &
  COR June 2024 arrivals, no WC adjustment &
  Lower bound; isolates baseline risk from routine June tourism \\[4pt]
Model 2: WC-adjusted &
  June 2026 projections via NTTO $\phi_c$ factors &
  Captures WC surge; $\phi_\text{WC}=1.174$ for qualified nations,
  $\phi_\text{base}=1.077$ for others \\[4pt]
Model 3: Schedule-driven &
  Travel decomposed into WC-fan stream (match schedule) +
  background stream (T-100) &
  Spatial routing of fans to the specific cities where their team plays \\
\bottomrule
\end{tabular}
\caption{Three-model hierarchy. All models use the same T-100 routing
fractions and disease parameters; they differ only in how travel
volumes are specified.}
\label{tab:models}
\end{table}

\subsection*{Data sources and parameters}

Travel volumes were derived from COR~I-94 monthly arrivals by country
of residence~\cite{CBP_I94}, scaled to June~2026 using NTTO
international visitor forecasts~\cite{NTTO2025} for the top-12
source markets, and to BTS T-100 International Segment data for
city-level routing fractions~\cite{BTS_T100}. Match schedules were
obtained from official FIFA records~\cite{FIFA2026_schedule}.
Disease incidence was derived from WHO surveillance data: dengue
from the WHO/PAHO global dengue dashboard~\cite{WHO_dengue},
malaria from the World Malaria Report~\cite{whoWMR2024}, measles
and pertussis from the WHO Immunization Data
portal~\cite{WHO_immunization}, and influenza from the WHO FluNet /
GISRS global sentinel surveillance system (ISO weeks 22--26,
2023--2025). Under-reporting factors and travel-while-infectious
probabilities were taken from the
literature~\cite{bhatt2013,undurraga2013,whoWMR2024,simons2012,mclaughlin2016}.
Key parameter values are: dengue ($\rho=0.10$, $p_d=0.50$), malaria
($\rho=0.20$, $p_d=0.30$), measles ($\rho=0.60$, $p_d=0.05$),
pertussis ($\rho=0.10$, $p_d=0.70$), and influenza ($\rho=0.10$,
$p_d=0.50$). Full details are given in the Supplementary Material.

\subsection*{Monte Carlo uncertainty propagation}

Both $\rho_d$ (under-reporting factor) and $p_d$ (travel-while-infectious probability) are uncertain across one or more orders of magnitude in the literature (Table~S2 in the Supplementary Material). To propagate this uncertainty through the importation estimates, we drew 5,000 independent samples of each parameter from independent Uniform distributions calibrated to the literature ranges: $\rho_d \sim \text{Uniform}(\rho_\text{min}, \rho_\text{max})$ and $p_d \sim \text{Uniform}(p_\text{min}, p_\text{max})$. Uniform distributions were chosen for transparency; they make no assumption about the shape of uncertainty within the documented range.

Because $\Lambda_{v,d} = p_d \cdot \rho_d \cdot \sum_c (C_{c,d}/P_c) \cdot N_{c,v}$ is linear in the product $(p_d \cdot \rho_d)$, each Monte Carlo draw rescales the central estimate by a factor $(p_d^{(i)} \cdot \rho_d^{(i)}) / (p_c \cdot \rho_c)$. We report the median and 95\% interval (2.5th--97.5th percentile) across the 5,000 draws. Figures show bars at the median with whiskers spanning the 95\% interval; the probability of at least one importation $P(\geq 1)$ is derived from the sampled $\Lambda$ distributions. Literature ranges used are: dengue ($\rho$: 0.06--0.26; $p$: 0.30--0.70), malaria ($\rho$: 0.10--0.35; $p$: 0.10--0.50), measles ($\rho$: 0.40--0.80; $p$: 0.02--0.10), pertussis ($\rho$: 0.01--0.10; $p$: 0.50--0.90), and influenza ($\rho$: 0.033--0.20; $p$: 0.30--0.70).

% ============================================================
\section*{Results}
% ============================================================

\subsection*{Importation risk overview}

Figure~\ref{fig:risk_heatmap} shows the probability of at least one importation
$[P(\geq 1)]$ for all five diseases across the 11~US venue cities under the
schedule-driven model (Model~3). Dengue fever leads importation risk at most
venue cities, with $P(\geq 1) > 0.99$ at Miami, New York, and Houston and
ranging from 0.09 to 0.97 at the remaining cities. Atlanta is a notable
exception: malaria reaches $P(\geq 1) > 0.99$---exceeding dengue
(0.97)---driven by direct travel from West and Central African nations.
Influenza ranks second at most cities, its risk amplified by coincidence with
the Southern Hemisphere winter peak. Pertussis shows broad but moderate
probability; outside Atlanta and New York, malaria risk is low and
geographically concentrated. Measles consistently shows the lowest probability,
driven by its severe prodrome suppressing
travel during the infectious window ($p_d = 0.05$). These five diseases
occupy distinct niches---by magnitude, geography, and uncertainty structure---as
the following subsections detail. The equivalent heatmaps for the Baseline
(Model~1) and WC-adjusted (Model~2) tiers are shown in the Supplementary
Material (Figures~S3--S4).

\begin{figure}[h!]
  \centering
  \includegraphics[width=\textwidth]{Figures/fig1_risk_heatmap.png}
  \caption{Importation risk overview for the FIFA World Cup 2026. Each cell
    shows $P(\geq 1\text{ importation})$ under the schedule-driven model
    (Model~3), June~2026, with Monte Carlo median from 5,000 parameter draws.
    Cities are ordered left to right by decreasing aggregate importation
    intensity; diseases are ordered top to bottom by decreasing overall
    probability. Dengue leads at most cities; malaria exceeds dengue at Atlanta
    ($P(\geq 1) > 0.99$ vs.\ 0.97). Measles risk is the lowest of the five
    diseases considered.}
  \label{fig:risk_heatmap}
\end{figure}

\subsection*{Importation intensity and parameter uncertainty}

Figure~\ref{fig:lambda_ci} shows the expected importation intensity ($\Lambda$,
median and 95\% Monte Carlo uncertainty interval) per city for each disease.
For dengue, New York leads (median $\Lambda = 12.4$, 95\%~CI 6.1--30.8),
followed by Los Angeles ($\Lambda = 11.8$, CI 5.8--29.2) and Miami
($\Lambda = 10.9$, CI 5.3--27.1), driven in all cases by Brazil's combination
of high endemic burden and high June travel volume. The dengue CI extends
substantially above the median because the central detection rate $\rho = 0.10$
lies near the lower end of the literature range (0.06--0.26): true importation
intensity may be considerably higher than the central estimate.

Pertussis shows the opposite asymmetry. Its central $\rho = 0.10$ is the
\emph{upper} bound of its literature range (0.01--0.10), so the 95\%~CI
extends far below the median. Pertussis estimates are best read as conservative
upper bounds: at the lower end of the plausible detection range, expected
importations could be as low as 7\% of the median. Malaria carries a moderate,
roughly symmetric CI. Measles and influenza uncertainties are modest relative
to their medians. The directional asymmetry of each CI is visualised in
Figure~S2 in the Supplementary Material: diseases with upward-skewed CIs
(dengue, influenza) may be underestimated; diseases with downward-skewed CIs
(pertussis) are conservative upper bounds.

\begin{figure}[h!]
  \centering
  \includegraphics[width=\textwidth]{Figures/fig2_lambda_ci.png}
  \caption{Expected importation intensity ($\Lambda$) with 95\% Monte Carlo
    uncertainty intervals under the schedule-driven model (Model~3).
    Each panel shows one disease; cities are ordered from highest to lowest
    median $\Lambda$. Points show the MC median; horizontal whiskers span the
    95\% interval (5,000 Uniform draws on $\rho_d$ and $p_d$).
    Note free x-axis scales: dengue intensity is orders of magnitude above
    the other diseases.}
  \label{fig:lambda_ci}
\end{figure}

The city and disease rankings visible in Figure~\ref{fig:lambda_ci} are robust
across all three model tiers (Figure~S5 in the Supplementary Material). The step from the Baseline to the WC-adjusted model produces a
near-uniform $\sim$17\% uplift in $\Lambda$ ($\phi_\text{WC}=1.174$; the WC
fan increment accounts for approximately 8.3\% of projected arrivals for
qualified nations); the further step to the schedule-driven model introduces
modest city-specific divergence without altering rankings. Background
tourism---travel that would occur regardless of the tournament---accounts for
the dominant share of importation risk even for strongly WC-associated nations
such as Brazil and Colombia; non-qualified countries contribute exclusively
through this background stream.

\subsection*{Source country drivers}

Figure~\ref{fig:country_drivers} shows the top~10 source countries by total
expected importations for each disease, coloured by world region, with 95\%
Monte Carlo uncertainty intervals. Country rankings are robust: the top three
contributors per disease are stable across the full parameter range. The
geographic pattern is disease-specific. Latin America---Brazil, Colombia,
Venezuela---dominates dengue and influenza, reflecting both high endemic burden
and the largest travel volumes to the US among WC-qualified nations. West and
Central Africa (Nigeria, Ghana, Senegal, Zaire/DRC) drives malaria risk.
European nations (Germany, the United Kingdom, France) lead pertussis and
measles, consistent with the documented pertussis resurgence and ongoing
measles circulation in Europe. Brazil is the single highest-risk source nation
overall, appearing in the top contributors for dengue, malaria, influenza, and
pertussis simultaneously.

\begin{figure}[h!]
  \centering
  \includegraphics[width=\textwidth]{Figures/fig4_country_drivers.png}
  \caption{Top~10 source countries by total expected importations
    (schedule-driven model, summed across all US venue cities) for each
    disease. Bars show the MC median; whiskers show the 95\% uncertainty
    interval. Colors indicate world region. Latin America dominates dengue
    and influenza risk; Africa drives malaria; Europe is the primary source
    for pertussis and measles.}
  \label{fig:country_drivers}
\end{figure}

% ============================================================
\section*{Extensions and Generalizability}
% ============================================================

Although applied here to the 2026 FIFA World Cup, the Poisson
importation framework is deliberately general. The three core
inputs---a travel volume matrix, a country-level disease incidence
table, and a travel-while-infectious probability---can be assembled
for any disease, any set of destination cities, and any international
gathering. The WC2026 application is a case study, not a constraint.

\textbf{Additional diseases.}
The model applies to any pathogen for which country-level incidence
and a travel-while-infectious probability can be specified. The
current analysis covers five diseases (dengue, influenza, malaria,
measles, pertussis); two near-term extensions are \textit{chikungunya}
and \textit{mumps}. Chikungunya poses particularly high importation
risk given the 2025--2026 outbreak in Brazil and Argentina (concurrent
analyses estimate $>$10 expected monthly imports from Brazil alone at
baseline~\cite{epistorm2026}); monthly case data are available from
PAHO and ECDC. Mumps data are available from the same WHO Immunization
portal used for measles and pertussis. Mpox (Clade I) and rubella
are lower-priority additions given low importation probabilities under
current outbreak conditions. For any newly emerging pathogen, the
framework requires only a case series by country and an estimate of
the travel-while-infectious window.

\textbf{Beyond venue cities.}
The current analysis covers the 11 US venue cities represented in
BTS T-100 routing data. The framework extends immediately to any
US city or county by substituting an alternative routing
source---DOT enplanement data, commercial booking records, or
domestic connection data---for the T-100 fraction $f_{c,v}$.
This is directly relevant for cities connected to venue cities by
domestic flights, where internationally arriving travelers may
complete their journeys on domestic legs not captured by T-100. The
same logic applies beyond the United States: any country with
international arrival records and disease surveillance data can be
modeled as a destination.

\textbf{Other mass gatherings and routine surveillance.}
The three-model hierarchy applies to any large international
gathering: Olympic Games, religious pilgrimages (e.g.,\ Hajj),
international summits, or music festivals. The only event-specific
inputs are the growth factor decomposition and the venue routing
schedule; all other components transfer directly. Beyond mass
gatherings, Model~1 functions as a standing importation surveillance
tool: updated monthly with current COR arrivals and WHO disease
data, it provides a near-real-time estimate of importation pressure
at any US airport city, independent of any scheduled event.

% ============================================================
\section*{Discussion}
% ============================================================

The dominant empirical finding of this analysis---that background tourism
accounts for the majority of importation risk even during a major
international sporting event---carries a direct practical implication:
public health preparedness should not focus exclusively on World Cup fans
but should treat the entire June international visitor population as the
primary exposure group.

The uncertainty analysis reveals important asymmetries in the reliability of
estimates across diseases. Dengue importation risk is robust across the full
parameter space, remaining near-certain at the highest-risk cities regardless
of detection-rate assumptions. Pertussis estimates, by contrast, should be
interpreted as conservative upper bounds: the central detection rate
($\rho = 0.10$) lies at the upper end of the literature range, so true
importation intensity is more likely to fall below the median than above it.
More broadly, diseases whose central detection rates fall near the lower bound
of the literature range (dengue, influenza) may be systematically
underestimated, a finding with direct implications for surveillance
prioritization.

The stability of city rankings across all three model tiers---with New
York, Los Angeles, Miami, and Houston consistently at highest aggregate
risk---demonstrates that the principal findings are robust to uncertainty
in WC-specific travel projections. A baseline analysis using only routine
June~2024 arrivals (Model~1) identifies the same priority cities as the
schedule-driven model (Model~3), supporting the practical utility of the
framework even when event-specific inputs are unavailable.

Quantifying influenza alongside vector-borne and vaccine-preventable
diseases reflects an important seasonal consideration: June coincides
with the Southern Hemisphere winter influenza peak, and WC-qualified
nations in South America and Oceania carry substantial seasonal burden
during the tournament window.

This framework is designed for reproducibility: all data
sources are publicly available and all code is openly shared, enabling
any public health team to deploy, update, and adapt the analysis without
proprietary data agreements. This positions the framework as
complementary to concurrent analyses employing airline booking records
and commercial mobility models~\cite{epistorm2026}, which are suited to
different institutional contexts. The WC2026 estimates are broadly
consistent with those of more resource-intensive models for endemically
stable diseases (dengue, malaria), validating the public-data approach
for first-order risk assessment. For outbreak-sensitive diseases,
discrepancies arising from the use of annual rather than rolling
surveillance data---most evident for UK measles and Norwegian
pertussis---quantify the premium on data currency for rapidly evolving
pathogens and motivate real-time data linkage as a priority
methodological improvement.

The three-model structure maps onto a practical planning timeline:
Model~1 can be executed as soon as host cities are announced;
Model~2 updated once attendance projections become available; and
Model~3 finalised once the match schedule is confirmed. City-specific
estimates and country-level contributions can inform decisions on which
venues and source-country streams to prioritise for enhanced
surveillance, though operational decisions require additional local
context---vector presence, population immunity, and healthcare
capacity---not captured here. The Monte Carlo uncertainty intervals
allow decision-makers to distinguish robust findings from estimates
warranting caution, supporting more calibrated risk communication than
a single point estimate permits. Beyond event-specific applications,
the simplicity of the data requirements opens the possibility of a
continuous monitoring mode: updated monthly with current COR arrivals
and WHO surveillance figures, the framework can track importation
pressure at US international gateways as a standing epidemiological
tool, independent of any scheduled mass gathering.

Several key limitations bound the scope of these estimates; full detail
is given in the Supplementary Material. BTS T-100 data cover nonstop
international segments only; passengers connecting through foreign hubs
lose their country identity, and passengers connecting through US hubs
are assigned to the gateway city rather than their final destination.
Disease incidence is derived from annual WHO reporting cycles, which
underestimate risk during active outbreaks (UK measles, European
pertussis surges) and may reflect prior-year dynamics rather than the
current transmission season. Growth factors $\phi_c$ are aggregate
projections; actual June~2026 travel volumes will differ. The model
estimates importation risk only---it does not model onward local
transmission, which depends on vector presence, population immunity,
and healthcare capacity not captured here. Finally, the Monte Carlo
analysis propagates uncertainty in disease parameters ($\rho_d$,
$p_d$) but treats travel volumes as fixed; uncertainty in the NTTO
projections, T-100 routing fractions, and the WC growth factor
decomposition would widen all intervals further.

The highest-priority methodological improvement is timelier disease
data: substituting WHO FluNet-style rolling surveillance averages for
annual incidence tables---where such data exist---would close the
data-currency gap most evident for measles and pertussis. Extensions
to additional diseases (chikungunya, mumps) and to domestic-connection
cities beyond T-100 gateways are described in the Extensions and
Generalizability section above. Longer term, coupling importation
estimates with local vector surveillance and population immunity data
would enable a full importation-to-transmission risk pathway, the
primary gap between what this framework currently provides and what
would be needed for outbreak probability estimation.

% ============================================================
\section*{Conclusions}
% ============================================================

We present a general, reproducible Poisson framework for quantifying
infectious disease importation risk that requires only publicly
available travel and surveillance data. The framework is applicable
to any disease, any destination city with international arrivals,
and any mass gathering or routine surveillance context. Applied to
the 2026 FIFA World Cup across five diseases and 11 US venue cities,
dengue fever emerges as the primary importation concern, with New
York, Los Angeles, Miami, and Houston at highest aggregate risk.
Influenza ranks second, reflecting its Southern Hemisphere winter
peak during the June tournament window. Background international
travel---not World Cup fan travel specifically---is the dominant
risk driver for all diseases, a finding with direct implications for
surveillance prioritization at any international event.

% ============================================================
\bibliographystyle{unsrtnat}
\bibliography{references}

\end{document}
